% ============================================================
%  6. Studija slučaja: Fiskalizacija 2.0
%  DOPRINOS RADA · ~27% opsega (cca 14 str. pri 52 str.; smije rasti)
%  Struktura 6.2 je zaključana; 6.1 i 6.3 slijede iz nje.
%
%  ⚠ TVRDO OGRANIČENJE: studija slučaja je REFERENTNA ARHITEKTURA
%    IZVEDENA IZ JAVNE SPECIFIKACIJE. Ni jedan stvarni sustav, poslodavac
%    ni klijent ne spominje se nigdje. Ni ranije studentski projekt.
%    Svaka tvrdnja mora imati citat propisa iza sebe.
%
% ============================================================
\chapter{Studija slučaja: Fiskalizacija 2.0}
\label{ch:studija}

\todo{Uvodni odlomak: što ovo poglavlje izvodi i zašto je izvod iz
specifikacije jači od opisa konkretnog sustava.}

% ------------------------------------------------------------
\section{Arhitektura i tijek dokumenta}
\label{sec:arhitektura}
% Budžet: ~5 str. (omjer, ne meta)
%
% SADRŽI:
%   - četverokutni model C1→C2→C3→C4 + CIS sa strane
%   - sudionici i njihove zakonske definicije [ZoF čl. 2.]
%   - DVA KANALA i DVA MODELA POUZDANOSTI:
%       razmjena PT↔PT: asinkrona, AS4 One-way/Push, maxretries=10,
%         potpisana potvrda primitka [PoeR čl. 6. st. 1. t. 8.]
%       fiskalizacija PT↔CIS: sinkroni SOAP [TS pogl. 8.],
%         bez timeouta, bez retryja, bez potvrde s pravnim učinkom
%     → NIJEDNA ODREDBA IH NE PREMOŠĆUJE
%   - granica povjerenja: C1↔C2, propis je ne uređuje uopće
%   - granica transakcije: između koraka 7 i 8 (dvije neovisne
%     fiskalizacije istog dokumenta, različiti rokovi)
%
% ✔ RIJEŠENO 2026-08-12 provjerom na primarnim dokumentima. Točno je:
%   - hrvatski profil: "Tehnički standardi i specifikacija povezivanja —
%     Pristupna točka i standardni AS4 profil", v1.4, 7.10.2025. (dok. 140)
%   - profil se zove eRačun-AS4 i sjedi na CEF eDelivery AS4 Profile
%     v1.15 (NE 1.14, NE Peppol AS4)
%   - "Peppol" se u dokumentu ne pojavljuje NIJEDNOM
%   - otkrivanje: AMS (Adresar Metapodatkovnih Servisa) + MPS
%     (Metapodatkovni servis), DNS/U-NAPTR → AMS, pa REST → MPS.
%     Peppol SML/SMP se NE koriste.
%   - Peppol legitimno ulazi samo kroz (a) kodne liste za prošireni
%     ISO 6523 (9934 = OIB, 0088 = GLN) i (b) naslijeđena imena shema
%     (busdox-docid-qns, cenbii-procid-ubl, iso6523-actorid-upis),
%     koja nose HRVATSKE vrijednosti.
%   - dva dodatna primarna dokumenta za sloj otkrivanja:
%     AMS v1.4, 9.10.2025 (dok. 142) i MPS v1.3, 8.10.2025 (dok. 141)
%
% ⚠ NE PISATI "BDXL 1.61" — ta verzija NE POSTOJI. Broj je artefakt
%   izvlačenja teksta iz PDF-a (oznaka fusnote "1" zalijepljena na "1.6").
%   Ispravno je BDXL 1.6 (30.5.2018.).
%
% SLIKA: dijagram tijeka u TikZ-u → rad/slike/

\todo{Model sudionika i tijek u jedanaest koraka; dva kanala i dva modela
pouzdanosti; granica povjerenja i granica transakcije. Dijagram.}

% ------------------------------------------------------------
\section{Identifikacija načina otkazivanja}
\label{sec:failure-modes}
% Budžet: ~7 str. — JEZGRA RADA. Smije rasti ako materijal nosi.
%
% OBLIK ARGUMENTA, BEZ IZNIMKE:
%   propis zahtijeva X [citat] → slijedi stanje S → S je neusklađen jer Z
% Iskustvo bira KOJE uzimamo ozbiljno; propis dokazuje da su MOGUĆI.
%
% NOSIVI CITAT CIJELOG POGLAVLJA — [ZoF čl. 48. st. 3.]:
%   "Fiskalizaciju eRačuna ... dužan je provoditi ODVOJENO OD POSTUPKA
%    RAZMJENE eRačuna u trenutku izdavanja eRačuna"
%   → zakon NALAŽE razdvajanje koje stvara problem konzistentnosti;
%     nesklad nije nedostatak implementacije nego PROPISANO STANJE.
%
% 12 RAZRAĐENIH SLUČAJEVA, GRUPIRANO UZROČNO (ne kao popis):
%   A  propisano razdvajanje kanala      A1, A2
%   B  sinkroni kanal bez semantike kvara B1, B2★, B3, B4
%      (B3: kardinalnost [1,100] s jednom greškom vrijedi za SVE ČETIRI
%       poslovne poruke, ne samo eRačun — sustavno svojstvo API-ja)
%   C  rokovi od nedefiniranih događaja   C1, C2, C3
%   D  ispad infrastrukture bez ventila   D1, D2
%   E  neregulirana granica C1↔C2         E1
%
% ★ B2 je najbolji pojedinačni izvod u radu. FORMULACIJA ISPRAVLJENA
%   2026-08-13 nakon čitanja primarnog dokumenta:
%   Identifikator eRačuna je SLOŽEN iz četiri polja, svako izrijekom
%   označeno u [TS pogl. 3.] kao "dio identifikatora eRačuna":
%     brojDokumenta (BT-1) + datumIzdavanja (BT-2)
%     + vrstaDokumenta (BT-3) + oibIzdavatelja (BT-31)
%   S008 provjerava jedinstvenost po tome + vrstaERacuna.
%   → KLJUČ ZA DEDUPLIKACIJU POSTOJI. Ne pisati da ga propis "uklanja".
%
%   Rupa je drugdje: [HRCIUS v1.8] nalaže da se ispravak bez poreznog
%   učinka šalje POD ISTIM BT-1 uz CopyIndicator="true", a CopyIndicator
%   NIJE dio identifikatora i k tome je 0..1. Specifikacija ne kaže
%   zadržava li ispravak isti BT-2:
%     - zadrži ga  → isti složeni identifikator → S008 (sudar s
%                    propisanim tijekom ispravka)
%     - dobije nov → nema sudara, ali ispravak ima drugi identifikator
%                    od originala, a u poruci NEMA polja koje bi original
%                    referenciralo (S17) → S012 nema na što
%   → S008 ISTOVREMENO znači "tvoj ponovni pokušaj je prošao" i "tvoj
%     propisani ispravak je odbijen", a ništa u poruci to ne razlikuje.
%
% KLASIFIKACIJA PO POSLJEDICI — četiri klase, od kojih je
%   "NEODREĐENO" metodološki najvažnija: nije da je ishod loš, nego da
%   ga sustav NE MOŽE SAZNATI. To je ono što obrasci mogu (najviše) suziti.
%   "Zakašnjelo bez fikcije urednosti" (C3) je VLASTITA KLASA — [čl. 21.
%   st. 3.] daje fikciju za B2C, [čl. 49.] nema ekvivalent, a [čl. 73.
%   st. 1. t. 21.] sankcionira.
%
% ★★ TEZA ODJELJKA — dvije vrste nesvodivosti:
%   FIZIKALNA (poznata literaturi): B1, B4 — dva generala, detekcija otkaza
%   REGULATORNO INDUCIRANA (klasa koju literatura ne opisuje):
%     A1/A2 — nesvodivi jer ih propis NALAŽE (čl. 48. st. 3. zabranjuje
%             zatvaranje jaza transakcijom)
%     B2    — nesvodiv jer propis nalaže tijek koji se sudara s vlastitom
%             provjerom duplikata, a ne daje način da se sudar razriješi
%     C3    — nesvodiv jer propis ne deklarira učinak sanacije
%     D1    — nesvodiv za obveznika (kvar države bez produljenja roka)
%
% PRILOG: preostalih 9 šutnji (S12–S14, S17, S19–S21, S23) kao tablica
%   s citatima, spomenuta jednom rečenicom. S22 ide u pogl. 9.

\todo{Izvod dvanaest načina otkazivanja u pet uzročnih skupina;
klasifikacija po posljedici; teza o dvjema vrstama nesvodivosti.}

% ------------------------------------------------------------
\section{Dizajn otpornog rješenja}
\label{sec:dizajn}
% Budžet: ~4 str.
% Bira iz rječnika poglavlja 4 prema načinima otkazivanja iz 6.2.
%
% NOSIVI IZVORI:
%   Helland 2007 — "entiteti" s vlastitim transakcijskim opsegom +
%     idempotentne at-least-once poruke među njima.
%     TO JE ODGOVOR NA A1/A2: kad transakcija preko kanala nije dopuštena,
%     entiteti + idempotentne poruke su ono što preostaje.
%   Schneider 1990 — stroj stanja kao odgovor na klasu "NEODREĐENO":
%     eksplicitno stanje "poslano-ali-nepotvrđeno" koje propis ne imenuje.
%   Saltzer i dr. 1984 — deduplikacija MORA biti na krajevima, ne u
%     transportu. Potkrijepljeno normativno: [OASISAS4 §3.4] kaže da se
%     semantika dostave ne smije pretpostaviti i da dedup nije obvezan.
%
% MORA ODGOVORITI, jer 6.2 to postavlja:
%   - kako razlikovati S008-kao-uspjeh od S008-kao-odbijeni-ispravak (B2)?
%   - kako se modelira stanje koje propis ne imenuje (B1)?
%   - što ERP smije zaključiti iz S008 nakon timeouta (B2, S6)?
%   - što se uopće može učiniti za regulatorno nesvodive slučajeve —
%     odgovor je vjerojatno DETEKCIJA I EVIDENCIJA, ne uklanjanje
%
% OTVORENO PITANJE IZ POGL. 5.2, ODGOVARA GA TIKET 12:
%   ima li fiskalizacija uopće što kompenzirati? Ako je jednosmjerna
%   zakonska obveza bez rollbacka, SAGA je tema pogl. 5 a ne uzorak ovdje.
%   To je legitiman nalaz, ali mora biti odlučen i obrazložen.

\todo{Izbor obrazaca prema izvedenim načinima otkazivanja; što se radi s
regulatorno nesvodivima; odgovor na pitanje o kompenzaciji.}
